\documentclass[../../main.tex]{subfiles}% 注意这里的文件路径不能用 ./main.tex ,否则用latexmk编译子文件会报错
\graphicspath{{\subfix{./image/}}{\subfix{../../image/}}} % 指定图片引用路径,后续可以直接使用图片文件名
% 如果图片存放在上级目录,则使用 ../../image/ 作为备用路径

% 例如:
% \begin{figure}[H]
% \centering
% \includegraphics[width=0.3\textwidth]{图.png}
% \caption{}
% \label{figure:图}
% \end{figure}
% 注意:上述\label{}一定要放在\caption{}之后,否则引用图片序号会只会显示??.

\begin{document}

\section{Cauchy积分公式的一些重要推论}

\begin{theorem}[Cauchy不等式]\label{theorem:Cauchy不等式(离散版)-复变函数}
设 \( f \) 在 \( B(a,R) \) 中全纯，且对任意 \( z \in B(a,R) \)，有 \( |f(z)| \leqslant M \)，那么
\begin{align}
|f^{(n)}(a)| &\leqslant \frac{n!M}{R^n}, \quad n = 1,2,\cdots. \label{eq:3.5.1}
\end{align}
\end{theorem}
\begin{note}
这个不等式给出了圆盘上全纯函数的各阶导数在圆心处值的估计.
\end{note}
\begin{proof}
取 \( 0 < r < R \)，则 \( f \) 在闭圆盘 \( \overline{B(a,r)} \) 中全纯，由\refthe{theorem:定理3.4.3}，得
\[
f^{(n)}(a) = \frac{n!}{2\pi \mathrm{i}} \int_{|\zeta - a| = r} \frac{f(\zeta)}{(\zeta - a)^{n + 1}} \mathrm{d}\zeta.
\]
于是，由\hyperref[theorem:长大不等式]{长大不等式}得
\[
|f^{(n)}(a)| \leqslant \frac{n!}{2\pi} \cdot \frac{M}{r^{n + 1}} \cdot 2\pi r = \frac{n!M}{r^n}.
\]
让 \( r \to R \)，即得所要证的不等式 \eqref{eq:3.5.1}。
\end{proof}

\begin{theorem}[Liouville定理]\label{theorem:Liouville(刘维尔)定理}
有界整函数必为常数。
\end{theorem}
\begin{proof}
设 \( f \) 为一有界整函数，其模的上界设为 \( M \)，即对任意 \( z \in \mathbb{C} \)，有 \( |f(z)| \leqslant M \)。任取 \( a \in \mathbb{C} \)，以 \( a \) 为中心、\( R \) 为半径作圆，因为 \( f \) 为整函数，故由 \hyperref[theorem:Cauchy不等式(离散版)-复变函数]{Cauchy 不等式}可得
\[
|f'(a)| \leqslant \frac{M}{R}.
\]
这个不等式对任意 \( R > 0 \) 都成立，让 \( R \to \infty \)，即得 \( f'(a) = 0 \)。因为 \( a \) 是任意的，所以在全平面上有 \( f'(z) \equiv 0 \)，因而由\refpro{proposition:复变函数导数为0必是常函数}可知\( f \) 是常数。
\end{proof}

\begin{theorem}[代数学基本定理]\label{theorem:代数学基本定理}
任意复系数多项式
\[
P(z) = a_0 z^n + a_1 z^{n - 1} + \cdots + a_n, \quad a_0 \neq 0
\]
在 \( \mathbb{C} \) 中必有零点。
\end{theorem}
\begin{remark}
考虑到实系数多项式在实数域中未必有零点，这个定理给出了复数域的又一重要性质。
\end{remark}
\begin{proof}

{\color{blue}证法一(\hyperref[theorem:Liouville(刘维尔)定理]{Liouville 定理}):}
反证,如果 \( P(z) \) 在 \( \mathbb{C} \) 中没有零点，那么 \( f(z) = \frac{1}{P(z)} \) 是一个整函数。由\( \lim_{z \to \infty} P(z) = \infty \)知
\begin{align*}
\lim_{z\rightarrow \infty} f\left( z \right) =\lim_{z\rightarrow \infty} \frac{1}{P(z)}=0.
\end{align*}
故存在$R>0$,使当 \( |z| > R \) 时，有
\begin{align*}
|f(z)|\leqslant 1.
\end{align*}
又由$f$是整函数知$f\in C(B(0,R))$,故由\refthe{theorem:连续的复变函数在紧集上的性质}知当 \( |z| \leqslant R \) 时,\( f \) 是有界的.因而 \( f \) 是有界整函数。由\hyperref[theorem:Liouville(刘维尔)定理]{Liouville 定理}，\( f \) 应是一常数,显然矛盾!这就证明了 \( P \) 在 \( \mathbb{C} \) 中必有零点。

{\color{blue}证法二(\hyperref[theorem:最大模原理-定理4.5.1]{最大模原理}):}假设多项式 \(P(z) = a_n z^n + a_{n-1} z^{n-1} + \cdots + a_1 z + a_0\)(其中 \(n \geqslant 1\) 且 \(a_n \neq 0\))在复平面 \(\mathbb{C}\) 上无零点.
令 \(f(z) = \frac{1}{P(z)}\).由于 \(P(z)\) 无零点,故 \(f(z)\) 为整函数.又因为当 \(|z| \to \infty\) 时,\(|P(z)| \to \infty\),所以 \(|f(z)| \to 0\).于是存在 \(R > 0\),使得当\(|z|>R\)时,有
\begin{align*}
|f(z)|<|f(0)|.
\end{align*}
注意到\(|f(z)|\)在\(\overline{B(0,R)}\)上连续,故存在\(z_0\in \overline{B(0,R)}\),使得
\begin{align*}
|f(z_0)|=\underset{z\in \overline{B(0,R)}}{\mathrm{sup}}|f(z)|\geqslant |f(0)|>|f(z)|, \quad \forall |z|>R.
\end{align*}
从而
\begin{align*}
|f(z_0)|=\underset{z\in \mathbb{C}}{\mathrm{sup}}|f(z)|\geqslant |f(z)|, \quad \forall z\in \overline{B(0,R+1)}.
\end{align*}
显然\(z_0\in \overline{B(0,R)}\subsetneq \overline{B(0,R+1)}\),因此\(|f(z)|\)在\(B(0,R+1)\)的内部取到最大值,这与\hyperref[theorem:最大模原理-定理4.5.1]{最大模原理}矛盾!
\end{proof}

\begin{theorem}[Morera(莫雷拉)定理]\label{theorem:Morera定理}
如果 \( f \) 是区域 \( D \) 上的连续函数，且沿 \( D \) 内任一可求长闭曲线的积分为零，那么 \( f \) 在 \( D \) 上全纯。
\end{theorem}
\begin{remark}
这个定理是\hyperref[theorem:Cauchy-Goursat定理(Cauchy积分定理)]{Cauchy积分定理的逆定理}.
\end{remark}
\begin{proof}
由\refthe{theorem:定理3.3.2}，存在单值全纯函数\( F\)，使得 \( F'(z) = f(z) \) 在 \( D \) 中成立。由\refthe{theorem:定理3.4.4}，\( F' =f\) 是 \( D \) 中的全纯函数.
\end{proof}

\begin{example}
设$f(z)$是整函数，若$f(\mathbb{C}) \subseteq \mathbb{C} \setminus [0,1]$，用刘维尔定理证明$f(z)$是常值函数。
\end{example}
\begin{note}
也可利用\hyperref[theorem:Picard定理-定理5.2.6]{Picard定理}立即得到证明.
\end{note}
\begin{proof}
令$F(z) = 2f(z) - 1$，则$F$是整函数，且
\[
f(\mathbb{C}) \subseteq \mathbb{C} \setminus [0,1] \Longrightarrow F(\mathbb{C}) \subseteq \mathbb{C} \setminus [-1,1].
\]
再令$g(z) = F(z) + \sqrt{F^2(z) - 1}$，则$g$是整函数，且在$\mathbb{C}$上能分出两个单值解析分支
\begin{align*}
g_1(z) &= F(z) - \sqrt{F^2(z) - 1}, \quad
g_2(z) = F(z) + \sqrt{F^2(z) - 1},
\end{align*}
其中$\sqrt{F^2(z) - 1}$恒指使得$\sqrt{F^2(z) - 1} \geqslant 0$的那一个单值解析分支。于是由$F(\mathbb{C}) \subseteq \mathbb{C} \setminus [-1,1]$和\hyperref[theorem:茹科夫斯基变换的基本性质]{茹科夫斯基函数的性质}知
\[
|g_1(z)| < 1, \quad |g_2(z)| > 1.
\]
由\hyperref[theorem:Liouville(刘维尔)定理]{Liouville定理}知$g_1(z)$，$\frac{1}{g_2(z)}$恒为常数。注意到$g_1(z) = \frac{1}{g_2(z)}$（$\forall z \in \mathbb{C}$），故可设$g_1(z) \equiv C$，$g_2(z) \equiv \frac{1}{C}$。因此
\[
F(z) = \frac{g_1(z) + g_2(z)}{2} = \frac{C + \frac{1}{C}}{2}, \quad \forall z \in \mathbb{C}.
\]
即$F(z)$也恒为常数，进而$f(z)$是常值函数。
\end{proof}

















\end{document}